\section{The Flip Audit Chain}
\label{sec:audit}

\subsection{Manifest Record}

Every \flip{} run writes the following fields to the plan manifest (the
JSON audit record that accompanies the submitted plan):

\begin{verbatim}
{
  "search": "flip",
  "base_seed": <u64>,
  "flip_steps": 10000,
  "percentile": 0.0,
  "visited_count": <usize>,
  "selected_plan_rank": <usize>,
  "chain_seed_0": <u64>
}
\end{verbatim}

Here \texttt{visited\_count} is $|\mathrm{visited}|$ (the number of distinct
plans stored, equal to 1 plus the number of accepted steps), and
\texttt{selected\_plan\_rank} is $\lfloor p \cdot |\mathrm{visited}| \rfloor$
(the 0-indexed position of the selected plan in the EC-sorted visited list).
The \texttt{chain\_seed\_0} field records the derived seed for deterministic
reproduction.

\subsection{Verification Protocol}

A verifier who wishes to confirm that the submitted plan is the correct
output of the stated \flip{} run proceeds as follows:

\begin{enumerate}
  \item Re-run the \flip{} chain with the recorded \texttt{base\_seed},
        \texttt{flip\_steps}, and \texttt{percentile} on the same census-tract
        graph and starting plan (which must also be recorded in the manifest).
  \item Confirm that the re-run's \texttt{visited\_count} matches the manifest.
        A mismatch indicates either a different starting plan or a different
        adjacency graph was used.
  \item Retrieve the plan at index \texttt{selected\_plan\_rank} in the
        EC-sorted visited list.
  \item Compare this plan's district assignment hash to the submitted plan's
        hash.
        A match confirms the submitted plan is the correct \flip{} output.
\end{enumerate}

This protocol is stronger than seed-only recording in two ways.
First, it confirms not just that the chain was run from a given seed, but
that the specific plan selected (by $p$-percentile) matches.
If an operator tried to substitute a different plan at the same percentile,
the hash comparison would fail.
Second, \texttt{visited\_count} provides an independent check on the chain's
execution: if the acceptance rate differs substantially from the expected
$50$--$65\%$ range (implying fewer or more than $5{,}000$--$6{,}500$ accepted
steps out of $10{,}000$), the verifier is alerted to investigate whether
the tract adjacency graph or population data differs from the reference.

\subsection{Why Seed-Only Recording Is Insufficient}

Consider the following attack: an operator runs the \flip{} chain, observes
that the plan at $p = 0.5$ produces a favourable partisan outcome, and
reports $p = 0.5$ in the manifest.
Later, they claim they actually ran the chain with $p = 0.0$ and got the same
plan ``by coincidence.''
Without \texttt{selected\_plan\_rank}, the verifier cannot distinguish these
scenarios: both produce the same seed, and the seed alone does not determine
which percentile was selected.

With \texttt{selected\_plan\_rank} recorded, the verifier re-runs the chain,
inspects the plan at rank $\lfloor 0.0 \times |\mathrm{visited}| \rfloor = 0$
(the minimum EC plan), and compares it to the submitted plan.
If they differ, the reported percentile was wrong.
This closes the manipulation window.

\subsection{Integration with the SHA-256 Audit Chain}

The \texttt{chain\_seed\_0} field is derived deterministically from the
\texttt{base\_seed} via the SHA-256 schedule:
\[
  \texttt{chain\_seed\_0} = \mathrm{SHA256}\!\left(
    \texttt{"FLIP\_CHAIN\_"} \,\|\, 0^{\mathrm{le64}} \,\|\,
    \texttt{"\_"} \,\|\, \texttt{base\_seed}^{\mathrm{le64}}
  \right)[0{:}8].
\]
Recording \texttt{chain\_seed\_0} explicitly lets the verifier confirm the
seed derivation without re-implementing the SHA-256 schedule; they can simply
check that \texttt{chain\_seed\_0} matches the expected derivation from
\texttt{base\_seed}.

The \texttt{base\_seed} itself is part of the broader plan manifest SHA-256
chain~\citep{b11paper}, which links the \flip{} run to the plan selection
decision (which METIS seed produced the starting plan), the data fetch
(which census-year shapefile was used), and the adjacency graph (which
boundary-sharing rules were applied).
This end-to-end chain makes it impossible to selectively manipulate any single
step without breaking the chain.
